\documentclass[12pt, a4paper]{article}
\usepackage[T2A]{fontenc}
\usepackage[utf8]{inputenc}
\usepackage[russian]{babel}
\usepackage{amsmath, amssymb, amsthm}
\usepackage{geometry}
\usepackage{bm}
\usepackage{arydshln}
\geometry{left=2.5cm, right=2cm, top=2cm, bottom=2cm}

\theoremstyle{definition}
\newtheorem*{theory}{Теория}
\newtheorem*{geom}{Геометрическое пояснение}
\newtheorem*{alg}{Алгебраическое пояснение}
\newtheorem{theorem}{Теорема}
\newtheorem*{remark}{Замечание}

\begin{document}

% Единый красивый заголовок
\begin{center}
    \Large\bfseries Задача №2: Сравнение производительности \texttt{memset} и цикла \texttt{for} \\
\end{center}

\vspace{0.5cm}

\section{Введение}
В данной работе исследуется производительность различных способов инициализации большого непрерывного блока памяти заданным значением (в данном случае --- символом пробела). В частности, проводится сравнительный анализ трех подходов:
\begin{enumerate}
    \item Обычный побайтовый цикл \texttt{for}.
    \item Стандартная библиотечная функция \texttt{memset}.
    \item Ускоренный цикл \texttt{for}, заполняющий память блоками по 8 байт (с использованием типа \texttt{uint64\_t}).
\end{enumerate}
Тестирование проводится на массиве размерностью $1$ ГБ ($1024 \times 1024 \times 1024$ байт).

\section{Анализ методов инициализации}

\subsection{Обычный цикл \texttt{for}}
\begin{alg}
Этот метод осуществляет классическую последовательную побайтовую запись. В большинстве современных архитектур ЭВМ такая запись крайне неэффективна. Процессор и шина памяти оптимизированы для передачи данных более крупными блоками (например, целыми кэш-линиями по 64 байта). Это наиболее медленный вариант, так как на каждый записываемый байт тратится несколько процессорных тактов (инкремент счетчика цикла, проверка условия выхода, разыменование указателя).
\end{alg}

\subsection{Ускоренный цикл \texttt{for} (8-байтовые блоки)}
\begin{alg}
Данный подход представляет собой ручную имитацию низкоуровневых оптимизаций средствами языка C++. Вместо побайтовой записи используется запись 64-битных (8 байт) блоков через приведение указателя к типу \texttt{uint64\_t*}. Значение искомого символа дублируется 8 раз в шестнадцатеричной маске (паттерне) \texttt{0x2020202020202020ULL}. Количество физических обращений к памяти и инструкций управления циклом уменьшается ровно в 8 раз. Хвост массива (остаток от деления размера на 8) аккуратно обрабатывается классическим побайтовым способом.
\end{alg}

\subsection{Функция \texttt{memset}}
\begin{theory}
Функция \texttt{memset} реализована в стандартной библиотеке языка C. Ее исходный код глубоко оптимизирован под конкретную архитектуру процессора на уровне ассемблера. Современные реализации используют векторные расширения команд для параллельной записи огромных блоков данных (по 16, 32 или 64 байта) за один такт процессора.
\end{theory}
\begin{remark}
Кроме того, в \texttt{memset} часто применяются инструкции потоковой записи (non-temporal stores), которые позволяют сбрасывать данные напрямую в оперативную память, минуя и не загрязняя кэш процессора (L1/L2/L3), что критически важно при инициализации массивов гигабайтных размеров.
\end{remark}

\hrule\vspace{0.5cm}

\section{Результаты численного эксперимента}
Для проверки производительности программа была скомпилирована компилятором \texttt{g++} и запущена в среде Ubuntu. Получены следующие усредненные результаты процессорного времени выполнения для инициализации $1$ ГБ памяти:

\begin{table}[h]
    \centering
    \begin{tabular}{|l|c|}
        \hline
        \textbf{Метод инициализации} & \textbf{Время выполнения (секунды)} \\
        \hline
        Обычный побайтовый цикл \texttt{for} & $3.64062$ \\
        Ускоренный цикл \texttt{for} (8 байт) & $0.23437$ \\
        Библиотечная функция \texttt{memset} & $0.07812$ \\
        \hline
    \end{tabular}
    \vspace{0.2cm}
    \caption{Сравнение времени инициализации памяти ($1$ ГБ)}
\end{table}

\section{Выводы}
Результаты эксперимента наглядно демонстрируют колоссальную разницу в подходах к работе с памятью. Обычный побайтовый цикл ожидаемо оказался самым медленным ($\approx 3.64$ с), что связано с огромным количеством холостых итераций и неэффективным использованием пропускной способности шины памяти.

Ускоренный цикл с блочной записью по 64 бита ($8$ байт) показал отличный результат ($\approx 0.23$ с), ускорив работу примерно в \textbf{15 раз} по сравнению с наивным подходом. Это произошло за счет радикального снижения накладных расходов на инструкции цикла и эффективной утилизации 64-битных регистров процессора.

Однако абсолютным лидером стала стандартная функция \texttt{memset} ($\approx 0.078$ с), оказавшись еще в \textbf{3 раза быстрее} ручной 64-битной оптимизации и почти в \textbf{46 раз быстрее} обычного цикла. Такой выдающийся результат достигается исключительно за счет глубоких аппаратно-зависимых оптимизаций (векторизация и кэш-байпас), недоступных в чистом высокоуровневом C++ коде без прямого использования интринсиков.

Таким образом, в промышленном коде всегда следует отдавать предпочтение стандартным функциям работы с памятью, так как их реализация максимально эффективно использует аппаратные возможности целевой платформы.

\end{document}